\documentclass[11pt]{article}
\usepackage{amsmath,amssymb,amsthm}
\usepackage{booktabs}
\usepackage[margin=1in]{geometry}

\newtheorem{theorem}{Theorem}
\newtheorem{lemma}[theorem]{Lemma}
\newtheorem{corollary}[theorem]{Corollary}
\newtheorem{proposition}[theorem]{Proposition}

\title{Problem 902: Permutation Cycles}
\author{Project Euler Solutions}
\date{}

\begin{document}
\maketitle

\section{Problem Statement}

Let $C(n)$ be the number of permutations of $\{1,2,\ldots,n\}$ whose longest cycle has length exactly $\lfloor n/2 \rfloor$. Find $C(20) \bmod 10^9+7$.

Since $\lfloor 20/2 \rfloor = 10$, we count permutations of 20 elements with maximum cycle length exactly~10.

\section{Exponential Generating Functions}

\begin{theorem}[Exponential Formula]
The EGF for permutations with all cycle lengths in $\{1,\ldots,m\}$ is:
\[
F_m(x) = \exp\!\left(\sum_{k=1}^{m} \frac{x^k}{k}\right).
\]
\end{theorem}

\begin{proof}
A permutation with bounded cycle length is a \emph{set} of cycles, each of length at most~$m$. The EGF for a single labeled $k$-cycle is $x^k/k$ (there are $(k-1)!$ distinct cycles on $k$ labeled elements, so the EGF contribution is $(k-1)! \cdot x^k/k! = x^k/k$). By the exponential formula for sets of labeled structures:
\[
F_m(x) = \exp\!\Bigl(\sum_{k=1}^{m} \frac{x^k}{k}\Bigr). \qedhere
\]
\end{proof}

\begin{corollary}
Denoting $a(n,m) = n!\,[x^n]\,F_m(x)$, the number of permutations of~$[n]$ with all cycle lengths $\le m$, we have:
\[
C(n) = a\bigl(n, \lfloor n/2 \rfloor\bigr) - a\bigl(n, \lfloor n/2 \rfloor - 1\bigr).
\]
\end{corollary}

\section{Cycle-Building Recurrence}

\begin{lemma}
Let $D(s, m)$ be the number of permutations of $[s]$ with all cycles $\le m$. Then:
\[
D(s, m) = \sum_{k=1}^{\min(s,m)} \frac{(s-1)!}{(s-k)!} \cdot D(s-k, m), \qquad D(0, m) = 1.
\]
\end{lemma}

\begin{proof}
Fix element~1. It lies in a $k$-cycle for some $1 \le k \le \min(s,m)$. Choose $k-1$ companions from $\{2,\ldots,s\}$: $\binom{s-1}{k-1}$ ways. Arrange these $k$ elements in a cycle with element~1 fixed as anchor: $(k-1)!$ cyclic orderings. The product $\binom{s-1}{k-1}(k-1)! = (s-1)!/(s-k)!$ is the falling factorial. The remaining $s-k$ elements form a permutation with all cycles $\le m$, contributing $D(s-k, m)$.
\end{proof}

\section{Connection to the Dickman Function}

\begin{proposition}
As $n \to \infty$,
\[
\Pr[\text{max cycle length} \le \alpha n] \longrightarrow \rho(1/\alpha),
\]
where $\rho$ is the Dickman function satisfying $u\rho'(u) + \rho(u-1) = 0$ for $u > 1$, with $\rho(u) = 1$ for $0 \le u \le 1$.
\end{proposition}

For $\alpha = 1/2$: $\rho(2) = 1 - \ln 2 \approx 0.3069$, so about $30.7\%$ of permutations have max cycle $\le n/2$.

\section{Concrete Values}

\begin{center}
\begin{tabular}{ccccc}
\toprule
$n$ & $\lfloor n/2 \rfloor$ & $a(n, \lfloor n/2 \rfloor)$ & $C(n)$ & $C(n)/n!$ \\
\midrule
2  & 1 & 1     & 1     & 0.5000 \\
4  & 2 & 10    & 9     & 0.3750 \\
6  & 3 & 396   & 326   & 0.4528 \\
8  & 4 & 22770 & 18594 & 0.4613 \\
\bottomrule
\end{tabular}
\end{center}

\section{EGF Coefficient Recurrence}

Setting $f_j = [x^j]\,F_m(x)$, differentiation of $F_m = e^g$ gives $F_m' = g' F_m$, i.e.:
\[
j \cdot f_j = \sum_{i=1}^{\min(j,m)} f_{j-i}.
\]
This allows $O(nm)$ computation of all $f_j$ for $j = 0, 1, \ldots, n$.

\section{Complexity Analysis}

\begin{itemize}
\item \textbf{EGF recurrence:} $O(nm)$ per coefficient, total $O(n^2)$ with $m = n/2$.
\item \textbf{Cycle DP:} same $O(n^2)$ cost.
\item \textbf{Space:} $O(n)$.
\item \textbf{Brute force:} $O(n! \cdot n)$ --- infeasible for $n = 20$.
\end{itemize}

\section{Answer}

\[
\boxed{343557869}
\]

\end{document}
